\section*{Hoofdstuk \ref{ch:g3:sorting-living-things} --- Levende wezens indelen}

\begin{solution}{exo:g3:sorting-living-things:1}
Een \omterm{def:g3:sorting-living-things:criterion}{criterium} is de regel die je bij het indelen gebruikt. Goed: een
lichaamskenmerk dat je kunt waarnemen --- ``heeft \omterm{def:g1:animals-around-us:covering}{veren}'', ``heeft zes
poten''. Vermeden: waar het woont of wat het toevallig doet --- ``woont
bij water'', ``vliegt''.
\end{solution}

\begin{solution}{exo:g3:sorting-living-things:2}
\omterm{prop:g3:sorting-living-things:groups}{Insecten}: zes poten en een lichaam in drie delen. \omterm{prop:g3:sorting-living-things:groups}{Vissen}: vinnen en
\omterm{def:g1:animals-around-us:covering}{schubben}, levend in het water. \omterm{prop:g3:sorting-living-things:groups}{Vogels}: \omterm{def:g1:animals-around-us:covering}{veren}, vleugels en een snavel.
\omterm{prop:g3:sorting-living-things:groups}{Zoogdieren}: \omterm{def:g1:animals-around-us:covering}{vacht} of haren, en moeders die hun jongen \omterm{def:g2:animals-grow-up:born}{melk} geven.
\end{solution}

\begin{solution}{exo:g3:sorting-living-things:3}
Bij: insect. Goudvis: vis. Struisvogel: vogel. Koe: zoogdier. Kikker:
amfibie.
\end{solution}

\begin{solution}{exo:g3:sorting-living-things:4}
Ze heeft \omterm{def:g1:animals-around-us:covering}{vacht} en geeft haar jongen \omterm{def:g2:animals-grow-up:born}{melk}, en ze heeft geen \omterm{def:g1:animals-around-us:covering}{veren} en
geen snavel. Vliegen is iets wat ze doet, niet iets wat een vogel
bepaalt --- de kenmerken zeggen zoogdier.
\end{solution}

\begin{solution}{exo:g3:sorting-living-things:5}
Hij heeft geen \omterm{def:g1:animals-around-us:covering}{schubben} en niet het kieuwen-en-vinnenbouwplan van een
vis --- en hij geeft zijn kalf \omterm{def:g2:animals-grow-up:born}{melk}. Zwemmen is zijn gewoonte; zijn
kenmerken zeggen zoogdier.
\end{solution}

\begin{solution}{exo:g3:sorting-living-things:6}
Nee. Het \omterm{def:g3:sorting-living-things:criterion}{criterium} voor \omterm{prop:g3:sorting-living-things:groups}{insecten} is zes poten (en een lichaam in drie
delen); met acht poten zakt de spin daarvoor. De spin hoort bij een
andere groep, hoe ze er op het eerste gezicht ook uitziet.
\end{solution}

\begin{solution}{exo:g3:sorting-living-things:7}
Madeliefje: \omterm{def:g1:plants-around-us:parts}{bloemen} en zaden, zachte groene \omterm{def:g1:plants-around-us:parts}{stengel}. Mos: nooit
\omterm{def:g1:plants-around-us:parts}{bloemen}, geen echte \omterm{def:g1:plants-around-us:parts}{stengel}. Den: een boom (houtige stam), \omterm{def:g1:plants-around-us:parts}{bladeren} als
naalden, geen \omterm{def:g1:plants-around-us:parts}{bloemen} maar kegels met zaden. Rozenstruik: \omterm{def:g1:plants-around-us:parts}{bloemen} en
zaden, houtige \omterm{def:g1:plants-around-us:parts}{stengels}.
\end{solution}

\begin{solution}{exo:g3:sorting-living-things:8}
Open antwoord. In de meeste straten winnen de \omterm{prop:g3:sorting-living-things:groups}{insecten} met grote
voorsprong --- mieren, vliegen, bijen, kevers, vlinders --- met de
\omterm{prop:g3:sorting-living-things:groups}{vogels} als tweede.
\end{solution}

\begin{solution}{exo:g3:sorting-living-things:9}
Die van de vriendin. ``Huisdieren tegenover wild'' deelt in naar hoe
\omterm{def:g1:animals-around-us:animal}{dieren} met mensen samenleven, en dat kan veranderen --- een zwerfkat is
binnen en buiten hetzelfde \omterm{def:g1:animals-around-us:animal}{dier}. ``\omterm{prop:g3:sorting-living-things:groups}{Zoogdieren} tegenover \omterm{prop:g3:sorting-living-things:groups}{vogels}'' deelt
in naar lichaamskenmerken die een leven lang blijven: het soort
\omterm{def:g3:sorting-living-things:criterion}{criterium} van de bioloog.
\end{solution}

\begin{solution}{exo:g3:sorting-living-things:10}
Stap 1: bedekking --- \omterm{def:g1:animals-around-us:covering}{veren} (dicht opeengepakt), een snavel, twee
poten, vleugels (als flippers). Stap 2: de jongen komen uit eieren.
Stap 3: de kenmerken passen bij de \omterm{prop:g3:sorting-living-things:groups}{vogels}. Stap 4: de indrukken ``zwemt
als een vis, kan niet vliegen'' zijn het er niet mee eens --- en
verliezen. De pinguïn is een vogel.
\end{solution}
